\documentclass{article}

\usepackage[english]{babel}
\usepackage[letterpaper,top=2cm,bottom=2cm,left=3cm,right=3cm]{geometry}
\usepackage{amsmath,amssymb}
\setlength{\parindent}{0pt}

\title{DSC 190/291: Learning Theory through Formal Proof and Proof Presentation\\Assignment 1}
\date{UCSD, Spring 2026}

\begin{document}
\maketitle

\section{Overview}

This assignment has three technical parts plus an AI usage report:
\begin{itemize}
\item Part A: set up your individual repository.
\item Part B: solve a theory problem from Week 1.
\item Part C: implement the Perceptron algorithm and verify its mistake bound experimentally.
\item Part D: write a short report describing how you used AI during the assignment.
\end{itemize}

\section{Part A: Repository Setup (10 points)}

Create an individual Git repository for this course. This repository will hold all your assignment submissions for the quarter.

\begin{enumerate}
\item Create a public GitHub repository for this course.
\item Add a \texttt{README.md} with your name (or alias) and a one-line description.
\item Set up an AI workflow file in the root of your repository:
\begin{itemize}
\item If you use Claude Code, create a \texttt{CLAUDE.md}.
\item If you use Codex (or another agent), create an \texttt{AGENTS.md}.
\item You may include both if you use multiple tools.
\end{itemize}
\item Create a directory \texttt{hw1/} for this assignment.
\item Share the repository link with the instructor.
\end{enumerate}

\section{Part B: Theory Problem (40 points)}

Throughout this part, let
$$
X = [0,1], \qquad Y = \{-1,+1\},
$$
and define
$$
h_\theta(x) = \operatorname{sign}(x-\theta),
$$
with the convention that
$$
\operatorname{sign}(z) =
\begin{cases}
+1, & z \ge 0, \\
-1, & z < 0.
\end{cases}
$$

We study the following structural assumption.

\subsection*{$\Delta$-Separated Threshold-Realizable Sequences}

A sequence $\bigl((x_t,y_t)\bigr)_{t=1}^T$ is called $\Delta$-separated threshold-realizable if there exist $\theta^* \in [0,1]$ and $\Delta > 0$ such that, for every round $t$,
$$
y_t = h_{\theta^*}(x_t)
\qquad \text{and} \qquad
|x_t-\theta^*| \ge \Delta.
$$

\subsection*{Problem 1: From Continuous Thresholds to a Finite Class}

Design a finite grid $G \subseteq [0,1]$ whose size depends only on $\Delta$, and consider the associated finite threshold class
$$
H_G = \{h_\theta : \theta \in G\}.
$$

Prove that for every $\Delta$-separated threshold-realizable sequence, there exists some $\widetilde{\theta} \in G$ such that
$$
h_{\widetilde{\theta}}(x_t) = y_t
\qquad \text{for all } t=1,\dots,T.
$$

Then use the Halving theorem from lecture to derive a mistake bound of order
$$
O(\log(1/\Delta))
$$
for an explicit online learner.

Your answer should clearly state:
\begin{itemize}
\item your choice of grid $G$,
\item the size of $G$ as a function of $\Delta$,
\item the resulting mistake bound.
\end{itemize}

\subsection*{Problem 2: A Positive Margin from Separation}

View thresholds as linear predictors by choosing a feature map
$$
\phi : [0,1] \to \mathbb{R}^d
$$
and a unit vector $u^*$ depending on $\theta^*$.

Prove that every $\Delta$-separated threshold-realizable sequence is linearly separable with margin at least $c\Delta$ for some absolute constant $c>0$ under your representation, while
$$
\|\phi(x)\| \le R
\qquad \text{for all } x \in [0,1]
$$
for some absolute constant $R$.

Then use the Perceptron theorem from lecture to derive an explicit mistake bound of order
$$
O(1/\Delta^2).
$$

Your answer should clearly specify your chosen feature map, the margin lower bound, the norm bound, and the final mistake bound.

\subsection*{Problem 3: Comparison and Interpretation}

Explain why the continuous-threshold impossibility phenomenon from lecture does not contradict Problem 1.

Then compare the two mistake bounds obtained in Problems 1 and 2. Why do they scale differently with $\Delta$? What is each argument measuring about the problem?

\subsection*{Problem 4: Optional Strengthening --- Audit an AI Proof}

An AI assistant gives the following argument:

\begin{quote}
Because every example is at least $\Delta$ away from the true threshold, continuous thresholds effectively form a class of size $O(1/\Delta)$. Therefore any online learner, including Perceptron, must make at most $O(\log(1/\Delta))$ mistakes on every $\Delta$-separated threshold-realizable sequence.
\end{quote}

Identify at least two mathematical problems with this argument. Then write a corrected statement that is true and prove it.

\section{Part C: Perceptron --- Implementation and Experiments (35 points)}

In lecture, we proved that the Perceptron makes at most $R^2/\gamma^2$ mistakes on data that is linearly separable with margin $\gamma$ and bounded norm $\|x_t\| \le R$.

Implement the Perceptron algorithm in Python, design a data-generation procedure, and run experiments addressing the following questions:

\begin{enumerate}
\item How do you generate data where you know the margin $\gamma$ and the bound $R$?
\item How does the number of mistakes $M$ scale with $1/\gamma^2$? Plot this relationship.
\item Is the theoretical bound $R^2/\gamma^2$ tight, or does the Perceptron do better in practice?
\item The bound is independent of the dimension $d$. Is this what you observe?
\item What happens as $\gamma \to 0$? Connect your observations to the threshold counterexample from lecture.
\item If you used AI to help write your code, how did you verify that the implementation is correct? What tests or checks would catch a subtle bug in the update rule or data generator?
\end{enumerate}

Include your code, plots, and a short discussion of your findings in \texttt{hw1/}.

\section{Part D: AI Usage Report (15 points)}

Write a short report, about half a page, describing how you used AI in this assignment. Address:

\begin{enumerate}
\item Which parts of the assignment did you use AI for?
\item Give at least one concrete example of an AI suggestion you accepted and why.
\item Give at least one concrete example of an AI suggestion you rejected or modified, and why.
\item How did you verify that the code and results were correct?
\end{enumerate}

Place this report in \texttt{hw1/} as a PDF or Markdown file.

\end{document}
